% GENERATED FILE. Edit canonical lessons, then run npm run generate:books.
% canonical-source-hash: fnv1a64:90f5f4b9cb16d68b
% canonical-lessons: MR-C364-samstha, MR-C364-samiti, MR-C364-sabha, MR-C364-pravakta, MR-C364-karyakarta

\chapter{Organizations}
\label{ch:mr-organizations}
\hlchaptermodality{364}

\begin{chapteropening}
\textbf{By the end of this chapter:} \emph{I can talk about institutions, committees and the people in them.}

Complete the last lesson of chapter 364: I can talk about institutions, committees and the people in them.
\end{chapteropening}

\section[\texorpdfstring{\mr{संस्था} (saṁsthā)}{संस्था (saṁsthā)}]{\mr{संस्था} (saṁsthā) — an institution; an organization}

\label{lesson:MR-C364-samstha}



\begin{quote}
Before the new one: say the Marathi for self-interest, then the Marathi for to incite.
\end{quote}

\subsection*{You'll want to know: \mr{संस्था}}
\textbf{\mr{संस्था}} — \emph{saṁsthā} — \textquotedblleft{}an institution; an organization\textquotedblright{}.

\begin{quote}
\mr{ही} \mr{संस्था} \mr{गरीब} \mr{मुलांना} \mr{शिकवते}.
\end{quote}

\begin{quote}
\emph{hī saṁsthā garīb mulānnā śikavte.}
\end{quote}

\begin{quote}
This organization teaches poor children.
\end{quote}

Feminine.

Feminine: \mr{ही} \mr{संस्था}; \mr{शिक्षणसंस्था}, an educational institution.

\begin{cousinweb}
\textbf{\mr{संस्था}} (\emph{saṁsthā}) is a tatsama, taken from Sanskrit \emph{saṁsthā}, \textquotedblleft{}standing together, an establishment\textquotedblright{} (Monier-Williams).
\end{cousinweb}

\begin{grammarlens}[title={What you've built}]
One more word for giving reasons and opinions.
\end{grammarlens}

\subsection*{Your turn}
\begin{itemize}
\raggedright
  \item \emph{Say it:} \emph{saṁsthā}
  \item \emph{Say it:} \emph{saṁsthā}, once more
  \item \emph{Say it:} \emph{hī saṁsthā garīb mulānnā śikavte.} — This organization teaches poor children.
  \item \emph{Recall:} say the Marathi for self-interest, then the Marathi for to incite, then say \emph{saṁsthā} again
\end{itemize}

\begin{tcolorbox}[breakable,colback=teal!4,colframe=teal!35!black,title={Before you move on}]
What does \mr{संस्था} mean? (\textquotedblleft{}An institution; an organization\textquotedblright{}.) Say it once more.
\end{tcolorbox}
\section[\texorpdfstring{\mr{समिती} (samitī)}{समिती (samitī)}]{\mr{समिती} (samitī) — a committee}

\label{lesson:MR-C364-samiti}



\begin{quote}
Before the new one: say the Marathi for to incite, then the Marathi for an institution.
\end{quote}

\subsection*{You'll want to know: \mr{समिती}}
\textbf{\mr{समिती}} — \emph{samitī} — \textquotedblleft{}a committee\textquotedblright{}.

\begin{quote}
\mr{शाळेने} \mr{पाच} \mr{जणांची} \mr{समिती} \mr{नेमली}.
\end{quote}

\begin{quote}
\emph{śāḷene pāc jaṇāñcī samitī nemlī.}
\end{quote}

\begin{quote}
The school appointed a committee of five.
\end{quote}

Feminine.

Feminine: \mr{समिती} \mr{नेमली}.

\begin{cousinweb}
\textbf{\mr{समिती}} (\emph{samitī}) is a tatsama, taken from Sanskrit \emph{samiti}, \textquotedblleft{}meeting, assembly, council\textquotedblright{} (Monier-Williams).
\end{cousinweb}

\begin{grammarlens}[title={What you've built}]
One more word for giving reasons and opinions.
\end{grammarlens}

\subsection*{Your turn}
\begin{itemize}
\raggedright
  \item \emph{Say it:} \emph{samitī}
  \item \emph{Say it:} \emph{samitī}, once more
  \item \emph{Say it:} \emph{śāḷene pāc jaṇāñcī samitī nemlī.} — The school appointed a committee of five.
  \item \emph{Recall:} say the Marathi for to incite, then the Marathi for an institution, then say \emph{samitī} again
\end{itemize}

\begin{tcolorbox}[breakable,colback=teal!4,colframe=teal!35!black,title={Before you move on}]
What does \mr{समिती} mean? (\textquotedblleft{}A committee\textquotedblright{}.) Say it once more.
\end{tcolorbox}
\section[\texorpdfstring{\mr{सभा} (sabhā)}{सभा (sabhā)}]{\mr{सभा} (sabhā) — an assembly; a public meeting}

\label{lesson:MR-C364-sabha}



\begin{quote}
Before the new one: say the Marathi for an institution, then the Marathi for a committee.
\end{quote}

\subsection*{You'll want to know: \mr{सभा}}
\textbf{\mr{सभा}} — \emph{sabhā} — \textquotedblleft{}an assembly; a public meeting\textquotedblright{}.

\begin{quote}
\mr{संध्याकाळी} \mr{गावात} \mr{सभा} \mr{भरली}.
\end{quote}

\begin{quote}
\emph{sandhyākāḷī gāvāt sabhā bharlī.}
\end{quote}

\begin{quote}
In the evening an assembly gathered in the village.
\end{quote}

Feminine.

\mr{सभा} \mr{भरणे}, an assembly meets; \mr{सभागृह}, a hall.

\begin{cousinweb}
\textbf{\mr{सभा}} (\emph{sabhā}) is a tatsama, taken from Sanskrit \emph{sabhā}, \textquotedblleft{}assembly, hall, council\textquotedblright{} (CDIAL 13167).
\end{cousinweb}

\begin{grammarlens}[title={What you've built}]
One more word for giving reasons and opinions.
\end{grammarlens}

\subsection*{Your turn}
\begin{itemize}
\raggedright
  \item \emph{Say it:} \emph{sabhā}
  \item \emph{Say it:} \emph{sabhā}, once more
  \item \emph{Say it:} \emph{sandhyākāḷī gāvāt sabhā bharlī.} — In the evening an assembly gathered in the village.
  \item \emph{Recall:} say the Marathi for an institution, then the Marathi for a committee, then say \emph{sabhā} again
\end{itemize}

\begin{tcolorbox}[breakable,colback=teal!4,colframe=teal!35!black,title={Before you move on}]
What does \mr{सभा} mean? (\textquotedblleft{}An assembly; a public meeting\textquotedblright{}.) Say it once more.
\end{tcolorbox}
\section[\texorpdfstring{\mr{प्रवक्ता} (pravaktā)}{प्रवक्ता (pravaktā)}]{\mr{प्रवक्ता} (pravaktā) — a spokesperson}

\label{lesson:MR-C364-pravakta}



\begin{quote}
Before the new one: say the Marathi for a committee, then the Marathi for an assembly.
\end{quote}

\subsection*{You'll want to know: \mr{प्रवक्ता}}
\textbf{\mr{प्रवक्ता}} — \emph{pravaktā} — \textquotedblleft{}a spokesperson\textquotedblright{}.

\begin{quote}
\mr{संस्थेच्या} \mr{प्रवक्त्याने} \mr{प्रश्नांची} \mr{उत्तरं} \mr{दिली}.
\end{quote}

\begin{quote}
\emph{saṁsthecyā pravaktyāne praśnāñcī uttaraṁ dilī.}
\end{quote}

\begin{quote}
The organization's spokesperson answered the questions.
\end{quote}

Masculine. Plural: \mr{प्रवक्ते}, \emph{pravakte}.

Plural \mr{प्रवक्ते}; oblique \mr{प्रवक्त्या}-.

\begin{cousinweb}
\textbf{\mr{प्रवक्ता}} (\emph{pravaktā}) is a tatsama, taken from Sanskrit \emph{pravaktṛ}, \textquotedblleft{}one who announces, a speaker\textquotedblright{} (Monier-Williams). The Sanskrit nominative is \emph{pravaktā}.
\end{cousinweb}

\begin{grammarlens}[title={What you've built}]
One more word for giving reasons and opinions.
\end{grammarlens}

\subsection*{Your turn}
\begin{itemize}
\raggedright
  \item \emph{Say it:} \emph{pravaktā}
  \item \emph{Say it:} \emph{pravaktā}, once more
  \item \emph{Say it:} \emph{saṁsthecyā pravaktyāne praśnāñcī uttaraṁ dilī.} — The organization's spokesperson answered the questions.
  \item \emph{Recall:} say the Marathi for a committee, then the Marathi for an assembly, then say \emph{pravaktā} again
\end{itemize}

\begin{tcolorbox}[breakable,colback=teal!4,colframe=teal!35!black,title={Before you move on}]
What does \mr{प्रवक्ता} mean? (\textquotedblleft{}A spokesperson\textquotedblright{}.) Say it once more.
\end{tcolorbox}
\section[\texorpdfstring{\mr{कार्यकर्ता} (kāryakartā)}{कार्यकर्ता (kāryakartā)}]{\mr{कार्यकर्ता} (kāryakartā) — an activist; a worker for a cause}

\label{lesson:MR-C364-karyakarta}



\begin{quote}
Before the new one: say the Marathi for an assembly, then the Marathi for a spokesperson.
\end{quote}

\subsection*{You'll want to know: \mr{कार्यकर्ता}}
\textbf{\mr{कार्यकर्ता}} — \emph{kāryakartā} — \textquotedblleft{}an activist; a worker for a cause\textquotedblright{}.

\begin{quote}
\mr{कार्यकर्त्यांनी} \mr{गावात} \mr{झाडं} \mr{लावली}.
\end{quote}

\begin{quote}
\emph{kāryakartyānnī gāvāt jhāḍaṁ lāvlī.}
\end{quote}

\begin{quote}
The activists planted trees in the village.
\end{quote}

Masculine. Plural: \mr{कार्यकर्ते}, \emph{kāryakarte}.

Plural \mr{कार्यकर्ते}; oblique \mr{कार्यकर्त्या}-.

\begin{cousinweb}
\textbf{\mr{कार्यकर्ता}} (\emph{kāryakartā}) is a learned compound built from Sanskrit \emph{kārya}, \textquotedblleft{}what is to be done\textquotedblright{}, and \emph{kartṛ}, \textquotedblleft{}doer\textquotedblright{} (Monier-Williams).
\end{cousinweb}

\begin{grammarlens}[title={What you've built}]
One more word for giving reasons and opinions.
\end{grammarlens}

\subsection*{Your turn}
\begin{itemize}
\raggedright
  \item \emph{Say it:} \emph{kāryakartā}
  \item \emph{Say it:} \emph{kāryakartā}, once more
  \item \emph{Say it:} \emph{kāryakartyānnī gāvāt jhāḍaṁ lāvlī.} — The activists planted trees in the village.
  \item \emph{Recall:} say the Marathi for an assembly, then the Marathi for a spokesperson, then say \emph{kāryakartā} again
\end{itemize}

\begin{tcolorbox}[breakable,colback=teal!4,colframe=teal!35!black,title={Before you move on}]
What does \mr{कार्यकर्ता} mean? (\textquotedblleft{}An activist; a worker for a cause\textquotedblright{}.) Say it once more.
\end{tcolorbox}
